
\section{MOV-Algorithmus}
\begin{lemma}
    Sei $E$ eine elliptische Kurve über dem Körper $\mathbb{F}_q$ mit
    \[
        \operatorname{Spur}(\phi_q)
        = q+1-\#E(\mathbb{F}_q)=0.
    \]
    Sei $n\in\mathbb{Z}_{>0}$. Existiert ein Punkt
    $P\in E(\mathbb{F}_q)$ der Ordnung $n$, so gilt
    \[
        E[n]\subseteq E(\mathbb{F}_{q^2}).
    \]
\end{lemma}

Dadurch kann das diskrete Logarithmusproblem auf einer supersingulären elliptischen Kurve mit $\operatorname{Spur}(\phi_q)=0$ auf ein diskretes Logarithmusproblem in der multiplikativen Gruppe $\mathbb{F}_{q^2}^{\times}$ reduziert werden. Da hierfür lediglich eine
Körpererweiterung vom Grad $2$ benötigt wird, können effizientere Verfahren zur Berechnung diskreter Logarithmen in endlichen Körpern
angewendet werden.

Auch für supersinguläre Kurven mit $\operatorname{Spur}(\phi_q)\neq 0$ ist eine solche Reduktion möglich. Dabei kann ein rweiterungsgrad $m\in\{3,4,6\}$ erforderlich sein. Diese Erweiterungsgrade sind weiterhin klein und ermöglichen daher eine Beschleunigung der Berechnung des diskreten Logarithmus.


\section{Quantencomputer}

Modern Crypography (S.397)
Laufzeit, Schlüsselgrößen, Sicherheit

Es findet sich eine einfache erklärung zu Quantencomputing in \cite{buell24}, S.227ff.

Kurze Erklärung klassische Computer 

PC hat elektromagnetisch betätigte Schalter. Der physikalische Zustand eines solchen Schalters heitß offen oder geschlossen, wird mathematisch durch 0 oder 1 repräsentiert




Alles schön Visualisieren
5.3.4. Quantensichere SignaturverfahrenWegen der Bedrohung der klassischen Signaturverfahren durch die fortschreitende Entwicklung
hinreichend großer Quantencomputer sollte eine Migration zu quantensicheren Verfahren erfolgen. Dies betrifft insbesondere sensitive Anwendungen mit langen Migrationszeiten, beispielsweise
im Zusammenhang mit Public-Key-Infrastrukturen.
Diese Technische Richtlinie empfiehlt den Einsatz eines quantensicheren Signaturverfahrens
grundsätzlich nur in Kombination mit einem klassischen Signaturverfahren. Die Hybridisierung
sollte so erfolgen, dass das hybride Signaturverfahren sicher ist, solange mindestens eines der
Verfahren sicher ist. Eine natürliche und robuste Hybridisierung ist die Konkatenation von einer
quantensicheren mit einer klassischen Signatur, so dass die konkatenierte Signatur genau dann
als gültig anerkannt wird, wenn sämtliche einzelnen Signaturen gültig sind. Hierbei sollte darauf
geachtet werden, Schlüsselmaterial für hybride Signaturen dediziert nur für diesen Zweck zu
erzeugen und nicht etwa auch für nicht-hybride Signaturen zu verwenden.
Der Sicherheit von hashbasierten Signaturverfahren liegen nur komplexitätstheoretische Annahmen über kryptografische Hashfunktionen zugrunde. Daher können hashbasierte Signaturen
unter der Voraussetzung, dass die Implementierungssicherheit von zustandsbehafteten und zustandslosen hashbasierten Verfahren sorgfältig berücksichtigt wird, grundsätzlich auch alleine (das
heißt nicht hybrid) zum Einsatz kommen und als eine gute Methode für die Erstellung langfristig
sicherer Signaturen dienen. \cite{bundesamt26}, S.59


Insbesondere die Crypto-Szene setzt darauf, etwa die ECDSA-Signaturen bei Bitcoin und Ethereum.
Der Einfluss geht aber auch weit darüber hinaus: Die Reisepässe einger europäischer Staaten (darunter auch Deutschland) nutzen ECC etwa für den Zugriff auf den integrierten Chip.

Dabei müssen wir uns hauptsächlich die asymmetrische Kommunikation, zum Beispiel verschlüsselte E-Mails oder verschlüsselten Dateiaustausch, anschauen. Symmetrische Verschlüsselung, also das einfache Abspeichern von Daten auf einem Server, kann bereits heute mit dem klassischen Algorithmus AES-256 erfolgen. Dieser Algorithmus gilt als quantensicher, da bisher noch kein Quanten-Algorithmus entdeckt wurde, der diesen brechen kann. 

Kyber als Quantensicherer Algorithmus